% =============================================================================
% SECTION 1: MOTIVATION AND PROBLEM STATEMENT
% =============================================================================

\section{Motivation and Problem Statement}

Material takeoff---the process of extracting quantities from architectural plans---remains a critical bottleneck in construction estimation.
Estimators must infer assembly intent (e.g., header configurations, king stud packs, blocking patterns) from plan conventions and building codes.
Even with accurate BOMs, material waste remains high without stock-aware cut optimization.
Downstream stakeholders (framers, suppliers, project managers) require machine-readable outputs and detailed build instructions, not just raw quantities.

\subsection{Industry Challenges}

The construction industry faces persistent productivity gaps~\cite{mckinsey2017construction}:

\begin{itemize}
  \item \textbf{Time burden}: Professional estimators spend 50--80\% of their time on quantity takeoffs, limiting the number of bids they can process~\cite{rics2019takeoff}.
  A typical 2,000 sq.ft. residential frame takeoff requires 8--12 hours of manual work.

  \item \textbf{Error rates}: Manual measuring and counting leads to 10--15\% error rates in material estimates~\cite{deloitte2023outlook}, resulting in costly change orders, project delays, and strained supplier relationships.
  Common errors include miscounting studs for wall intersections, incorrect header sizing, and missing blocking requirements.

  \item \textbf{Inconsistency}: Two estimators working from the same plans often produce BOMs differing by 5--10\% in key materials like lumber and fasteners, undermining cost predictability.

  \item \textbf{Knowledge gaps}: Junior estimators lack deep knowledge of framing methodology and jurisdiction-specific code requirements (IRC~\cite{irc2021}, local amendments), leading to non-compliant or over-conservative designs.

  \item \textbf{Material waste}: Without cut optimization, lumber waste rates of 15--20\% are typical; optimized systems demonstrate waste reduction to under 5\%~\cite{manrique2007framex}, representing significant cost savings on large projects.
\end{itemize}

\subsection{Goals}

\textbf{Goal:} Create an end-to-end workflow that:
\begin{enumerate}
  \item[(a)] Interprets plans into \emph{component assemblies} (walls, headers, openings)
  \item[(b)] Selects \emph{real-world purchasable stock} (lengths, grades, treatments)
  \item[(c)] Computes \emph{cuts with waste minimization} using optimization algorithms
  \item[(d)] Exports data in \emph{supplier-friendly formats} (CSV, JSON, EDI)
  \item[(e)] Generates \emph{build instructions} with code compliance citations
\end{enumerate}
